%% kapa: exam=1 round=24 year=2013 period=1 form=A subject=civil_law subject_ko=민법 seq=10 number=10
%% source: Q-Net 감정평가사 기출문제 (cst00309) artlSeq=5073802 file=extracted/1008/5073802/2013년 제24회 감정평가사 1차시험 기출문제/1교시(A형).hwp
%% conversion: kordoc 4.24.51 -> tools/gen1.py (auto); review=proofread-sonnet; flags=none
\begin{question}{10}
사회질서에 반하는 법률행위에 관한 설명으로 옳지 않은 것은? (다툼이 있으면 판례에 의함)
\begin{choices}
  \item 부첩(夫妾)관계의 종료를 해제조건으로 하는 증여계약은 사회질서에 반하여 무효이다.
  \item 어떠한 일이 있어도 이혼하지 않겠다는 각서를 써 주었더라도 그와 같은 의사표시는 공서양속에 위반되어 무효이다.
  \item 법률행위가 사회질서에 위반되었는지의 여부는 법률행위의 효력발생시점을 기준으로 판단해야 한다.
  \item 강제집행을 면할 목적으로 허위의 근저당권설정등기를 하는 행위는 사회질서에 반하는 것으로 볼 수 없다.
  \item 소송에서 증언하는 대가로 통상적으로 용인될 수 있는 수준을 넘는 급부를 제공하기로 하는 약정은 사회질서에 반하는 것으로 무효이다.
\end{choices}
\end{question}
\answer{3}
